(1) 每个空间本征态 \(n\) 有两个自旋投影可用，而同种费米粒子不能处于完全相同的空间与自旋状态。因为 \(N=2K\) 为偶数，基态在 \(n=1,2,\ldots,K\) 每个空间轨道填入一对反向自旋粒子，\(n>K\) 均不占据，因此不需对最高壳层的简并部分填充作额外选择。
(2) 每个能级占两粒子，故 \(E_0=2\sum_{n=1}^K\varepsilon_n=(\pi^2\hbar^2/(mL^2))\sum_{n=1}^Kn^2\)。使用平方和恒等式 \(\sum_{n=1}^Kn^2=K(K+1)(2K+1)/6\)，得到精确 \(E_0=\pi^2\hbar^2K(K+1)(2K+1)/(6mL^2)\)。
(3) 费米能即最高已占空间轨道能量 \(E_F=\pi^2\hbar^2K^2/(2mL^2)\)。对 \(K\gg1\)，总基态能 \(E_0\sim\pi^2\hbar^2K^3/(3mL^2)\)，因此 \(E_0/(NE_F)\to1/3\)。这一系数是一维态密度积分的结果，与三维自由费米气体 \(3/5\) 的系数不同。
(4) 在固定粒子数下 \(E_0\propto L^{-2}\)，端点力 \(F=-\partial E_0/\partial L=2E_0/L>0\)，表示系统推开边界的作用。即使温度为零，填满的费米态仍携带动能并产生机械压力；其来源是边界量子化与泡利填充，而不是热理想气体的温度压强。
